\section{Experimentation}
\label{sec:exp}

This section outlines the experimental methodology and assesses the performance of NCAD-CS across both univariate spacecraft telemetry (NASA SMAP/MSL) and complex multivariate benchmarks from the mTSBench suite.

\subsection{Experimental Setup}
\label{subsec:exp_setup}
The experiments evaluate both univariate telemetry channels and multivariate cyber-physical systems (Daphnet, OPPORTUNITY, Exathlon, CalIt2, GECCO, Genesis, Room-Occupancy). The NCAD-CS v5 configuration employs the following settings:
\begin{itemize}[noitemsep, topsep=0pt, leftmargin=*]
    \item \textit{Windowing:} A sliding window scheme with context size $C=256$ and suspect/successor window $S=64$ (full window size $L=320$).
    \item \textit{Feature Representation:} Telemetry channels are standardized and augmented with multi-domain feature extraction and Takens time-delay attractor embeddings ($\tau=4, m=5$).
    \item \textit{Spatial-Temporal Relational Encoder:} The architecture uses the \texttt{RelationalGATEncoder}, fusing 4 causal dilated TCN layers with multi-head spatial Graph Attention Network (GAT) layers ($H=4$ heads) to dynamically model cross-variable dependencies alongside temporal trajectory and statistical-spectral moments.
    \item \textit{Successor Memory:} Counterfactual Successor Memory stores RPCA-sanitized $(h_{\text{ctx}}, h_{\text{succ}})$ pairs with $K=8$ nearest normal neighbors and an adaptive OOD gain mechanism.
    \item \textit{EVT Adaptive Tail Calibration:} In place of ad-hoc empirical elbow floors, the decision threshold is calibrated using Extreme Value Theory (EVT / SPOT). An asymptotic Generalized Pareto Distribution (GPD) is fitted to the tail of normal training residual excesses using Grimshaw's exact 1D root-finding algorithm to establish an extreme quantile threshold for risk level $q=10^{-3}$.
\end{itemize}

\subsection{Multi-Dataset Benchmark Evaluation}
\label{subsec:multidataset_results}
We evaluate NCAD-CS v5 across 23 diverse benchmark channels spanning 7 datasets from the mTSBench suite. Table~\ref{tab:multidataset_benchmark} reports the comparative performance of baseline NCAD against NCAD-CS v5 and the theoretical Oracle upper bound.

\begin{table}[h!]
\centering
\caption{Consolidated Multi-Dataset Benchmark Evaluation comparing Baseline NCAD, NCAD-CS v5 (Relational GAT + EVT Calibration), and Oracle Optimal PA-F1.}
\label{tab:multidataset_benchmark}
\begin{tabular}{@{}lcccc@{}}
\toprule
Dataset & Channels & Baseline NCAD & NCAD-CS v5 & Oracle Upper Bound \\
\midrule
Exathlon (Streaming Traces) & 6 & 0.4850 & \textbf{0.6692} & 0.7787 \\
Daphnet (Gait Accelerometry) & 6 & 0.3041 & \textbf{0.3517} & 0.6179 \\
Room-Occupancy (Environmental) & 2 & 0.1860 & \textbf{0.3721} & 0.7496 \\
OPPORTUNITY (77 Sensors) & 6 & 0.1253 & \textbf{0.1834} & 0.4922 \\
GECCO (Water Quality) & 1 & 0.0820 & \textbf{0.1928} & 0.5378 \\
CalIt2 (Traffic Flow) & 1 & 0.0702 & \textbf{0.0923} & 0.1380 \\
Genesis (Robotics) & 1 & 0.0610 & \textbf{0.0949} & 0.2342 \\
\midrule
\textbf{Macro Average} & \textbf{23} & \textbf{0.2604} & \textbf{0.3631 (+39.4\%)} & \textbf{0.5975 (+52.5\%)} \\
\bottomrule
\end{tabular}
\end{table}
